\chapter{Hasil dan Pembahasan}

% ================================================================
% PANDUAN PENULISAN BAB IV — HASIL DAN PEMBAHASAN
% Sajikan temuan secara objektif pada bagian Hasil, lalu tafsirkan dan
% kaitkan dengan teori serta penelitian lain pada bagian Pembahasan.
% ================================================================

\section{Hasil Penelitian}

% Deskripsikan karakteristik subjek dan hasil analisis sesuai dengan
% tujuan penelitian. Sajikan data berupa narasi yang didukung tabel
% dan/atau gambar. Bagian ini TIDAK boleh berisi penafsiran/interpretasi.

% Contoh gambar (ganti placeholder berikut dengan gambar Anda;
% simpan berkas gambar di folder figures/):
\begin{figure}[ht]
\centering
\fbox{\rule{0pt}{3cm}\rule{0.9\linewidth}{0pt}}
\caption{Tulis keterangan gambar di sini.}
\label{fig:contoh}
\end{figure}

% Alternatif bila gambar sudah tersedia:
% \begin{figure}[ht]
% \centering
% \includegraphics[width=0.8\linewidth]{figures/nama_file_gambar}
% \caption{Tulis keterangan gambar di sini.}
% \label{fig:contoh2}
% \end{figure}

% Contoh tabel:
\begin{table}[ht]
\centering
\begin{tabular}{ll}
\toprule
Variabel & Keterangan \\
\midrule
A & Contoh data A \\
B & Contoh data B \\
\bottomrule
\end{tabular}
\caption{Tulis keterangan tabel di sini.}
\label{tab:contoh}
\end{table}

\section{Pembahasan}

% Interpretasikan hasil penelitian dan kaitkan dengan teori serta hasil
% penelitian sebelumnya. Jelaskan temuan yang mendukung maupun yang
% bertentangan, serta berikan argumen yang logis atas perbedaan tersebut.

\section{Keterbatasan Penelitian}

% Akui kelemahan penelitian (desain, besar sampel, potensi bias) dan
% jelaskan implikasinya terhadap interpretasi serta generalisasi hasil.
